\documentclass[11pt,letterpaper]{article}
\usepackage{cornell-notes}
\usepackage{mathtools}

\title{Chapter 18: Three Short Topics}
\author{}
\date{}

\setDocTitle{Chapter 18: Three Short Topics}
\setDocSubtitle{String Algorithms}
\setDocVersion{v1.0}
\setDocStatus{Study Companion}
\setDocOwner{String Algorithms Cornell Notes}
\setDocAuthor{}
\setDocDate{\today}
\setDocSummary{Cornell-format study and review notes for Chapter 18: Three Short Topics.}

\setCornellCollection{String Algorithms}
\setCornellUnitType{Chapter}
\setCornellUnitNumber{18}
\setCornellUnitTitle{Three Short Topics}
\setCornellCompanionLabel{Study and review companion}

\hypersetup{pdftitle={Chapter 18: Three Short Topics},pdfsubject={Cornell notes},pdfauthor={}}

\begin{document}
\maketitle

\begin{CornellOverviewBox}{Chapter overview}
\textbf{Purpose.} Apply sequence-analysis methods to translated alignment with frameshifts, computational gene prediction, and molecular computation performed with DNA strands.

\textbf{Learning objectives}
\begin{itemize}
  \item Model reading-frame changes in DNA-to-protein matching.
  \item Combine coding signals and content evidence for gene prediction.
  \item Explain how DNA hybridization can implement combinatorial operations.
  \item Evaluate molecular parallelism without ignoring material complexity.
\end{itemize}

\textbf{Prerequisites.} Codons and translation, alignment DP, probabilistic sequence models, and graph decision problems.
\end{CornellOverviewBox}

\begin{CornellNotesTable}
\CornellNoteRow{\S18.1: Why are frameshifts special?}{DNA is translated in codons of three bases. An insertion or deletion not divisible by three changes the reading frame and alters every downstream codon until another compensating event. A direct character alignment does not capture this structured effect.}
\CornellNoteRow{How can DNA be matched to protein?}{Use DP states for the current DNA position, protein position, and reading-frame or codon phase. Transitions translate a codon to an amino acid, introduce amino-acid substitutions or gaps, or consume one or two extra bases with a frameshift penalty. State prevents a frame change from being treated as independent local mismatches.}
\CornellNoteRow{What is the frameshift profile?}{\textbf{Input/output:} DNA, protein, and a best translated alignment. \textbf{Model:} codon translation plus substitution, gap, and frameshift costs. \textbf{Invariant:} each state records the best score for its prefix pair and phase. \textbf{Checks:} stop codons, ambiguous bases, sequence orientation, and partial terminal codons.}
\CornellNoteRow{\S18.2: What evidence predicts a gene?}{Signals include start and stop codons, splice donor and acceptor sites, promoters, reading-frame consistency, and coding-region composition. No single signal is sufficient: short motifs occur by chance, while real genes vary across organisms.}
\CornellNoteRow{How are signals combined?}{Represent intergenic, exon, intron, and boundary states in a dynamic program or probabilistic state model. Transitions enforce legal gene structure; emissions or segment scores measure sequence evidence. The highest-scoring path provides predicted features and coordinates.}
\CornellNoteRow{What makes evaluation honest?}{Measure boundary accuracy as well as gene-level detection, separate training and test regions, account for alternative splicing and pseudogenes, and report organism-specific assumptions. A high aggregate accuracy can hide poor detection of rare boundaries.}
\CornellNoteRow{\S18.3: What is molecular computation?}{Encode candidate objects as DNA strands and use biochemical operations such as synthesis, hybridization, ligation, amplification, separation, and sequencing to create and filter large collections in parallel. The molecules are the data and the laboratory operations implement the computation.}
\CornellNoteRow{How can a path problem be encoded?}{Assign oligonucleotides to vertices and complementary overlaps to legal edges. Hybridization and ligation assemble strands representing walks; laboratory filters retain strands with the required length, endpoints, and vertex coverage. Detecting a surviving strand witnesses a solution under the encoding.}
\CornellNoteRow{Where is the hidden cost?}{Massive molecular parallelism may require exponentially many distinct molecules, careful error control, several sequential laboratory rounds, and expensive readout. Count material volume, error rates, and experimental steps rather than only the depth of the abstract operation sequence.}
\CornellNoteRow{\S18.4: What should practice emphasize?}{Add frame state to a small alignment recurrence, design a legal-state diagram for a simple gene, and audit a DNA-computing proposal for encoding uniqueness, false hybridization, molecule count, and verification cost.}
\end{CornellNotesTable}

\begin{CornellSummaryBox}{Chapter synthesis}
These topics show three roles for string models. Frameshift-aware DP aligns across two alphabets while tracking codon phase; gene prediction finds a highest-scoring path through structural states; DNA computing turns carefully designed strings into physical carriers of combinatorial candidates. In every case, the state representation is what makes the algorithm faithful to the domain.
\end{CornellSummaryBox}

\begin{CornellSummaryBox}{Key takeaways}
\begin{CornellRecallList}
  \item A frameshift has persistent phase effects.
  \item Translated alignment needs codon-aware state and scoring.
  \item Gene prediction combines boundary signals with region content.
  \item A legal-state path encodes a coherent gene structure.
  \item DNA computing performs operations on molecular encodings.
  \item Molecular parallelism can hide exponential material and error costs.
\end{CornellRecallList}
\end{CornellSummaryBox}

\begin{CornellOverviewBox}{Algorithm selection: when to use this}
\begin{CornellChecklist}
  \item Use phase-aware DP for DNA-to-protein comparison with possible frameshifts.
  \item Use a state model for gene prediction when signals and regions must form a legal structure.
  \item Use molecular computation only when its parallelism offsets encoding, material, and readout costs.
\end{CornellChecklist}
\end{CornellOverviewBox}

\begin{CornellExamBox}{Self-test questions}
\begin{CornellRecallList}
  \item Why does a one-base indel affect more than one translated symbol?
  \item Which states are needed for a simple frameshift-aware DP?
  \item Why are start and stop codons insufficient for gene prediction?
  \item How does a state path encode exon-intron structure?
  \item What laboratory operations implement selection in DNA computing?
  \item Why can shallow molecular computation still consume exponential resources?
\end{CornellRecallList}
\end{CornellExamBox}

{\footnotesize
\textbf{Terms and notation to remember.} codon, reading frame, frameshift, translated alignment, splice signal, exon, intron, state path, oligonucleotide, hybridization, ligation, molecular parallelism.

\textbf{Connections.} The same DP discipline used for alignment supports gene-structure paths, while careful string encodings turn abstract graph constraints into molecular compatibility.
}

\end{document}
